\section{Experiments}\label{sec:experiments}

For the resistor values we used E24 series.

\subsection{Inverting amplifier}\label{sec:inverting}

The first operational amplifier circuit to construct and measure was the
inverting amplifier. We have chosen the gain of $G=-10$ and the feedback
restistor value of $R_F=\qty{100}{\kilo\ohm}$. To calculate the value of $R_G$
resistor we transformed the~\autoref{eqn:inverting} into:
\begin{equation}
  R_G = -\frac{R_F}{G} = \frac{\qty{100}{\kilo\ohm}}{10} = \qty{10}{\kilo\ohm}
\end{equation}

Applying these values to our testbench PCB gave us $R_1=\qty{10}{\kilo\ohm}$,
$R_4=\qty{100}{\kilo\ohm}$. We calulated the compensation restistor $R_5$ value
as proposed in the laboratory instruction to get
$R_5=R_1|R_4=\frac{100}{11}\,\unit{\kilo\ohm} \approx \qty{9.1}{\kilo\ohm}$.
The circuit is shown in \autoref{fig:inverting}.

\begin{figure}[ht]
  \centering
  \input{schematics/inverting}
  \caption{Inverting amplifier without the feedback capacitor $C_3$}
  \label{fig:inverting}
\end{figure}

After measuring the circuit frequency response we added the feedback capacitor
$C_3=\qty{1}{\nano\farad}$ in parallel with $R_4$ as shown
in~\autoref{fig:inverting_c3}.

\begin{figure}[ht]
  \centering
  \input{schematics/inverting_c3}
  \caption{Inverting amplifier with the feedback capacitor $C_3$}
  \label{fig:inverting_c3}
\end{figure}

\todo[inline]{Design: chosen gain, $R_1$, $R_4$, $R_5 = R_1 \parallel R_4$, calculated $f_H$, $f_c$ with $C_3$ and worst-case offset.}

\subsubsection{Frequency response}
\todo[inline]{Bode plot (magnitude in dB and phase, logarithmic frequency): (a) without $C_3$, (b) with $C_3 = \qty{1}{nF}$, (c) cut-off frequencies marked and compared with the calculation.}

\subsubsection{SPICE analysis}
\begin{figure}[ht]
  \centering
  \includegraphics[width=\linewidth]{sim_inverting_ac}
  \caption{Inverting amplifier without $C_3$, AC analysis (ngspice, TL061 model)}
  \label{fig:sim_inverting}
\end{figure}
\begin{figure}[ht]
  \centering
  \includegraphics[width=\linewidth]{sim_inverting_c3_ac}
  \caption{Inverting amplifier with $C_3$, AC analysis (ngspice, TL061 model)}
  \label{fig:sim_inverting_c3}
\end{figure}
\todo[inline]{Compare the simulated cut-off frequencies (\autoref{fig:sim_inverting} and~\ref{fig:sim_inverting_c3}, \texttt{simulation/out/results.csv}) with the calculated ones and the measurement.}

\subsubsection{Offset}
\todo[inline]{Measured output offset with the input grounded, input-referred value, comparison with the data sheet.}

\subsubsection{Slew rate}
\todo[inline]{Oscilloscope screenshot of the output edge, measured slew rate, comparison with \qty{3.5}{V/\us} typical.}

\subsection{Non-inverting amplifier}\label{sec:noninverting}
\todo[inline]{Design: chosen gain, $R_1$, $R_4$, $R_2 = R_1 \parallel R_4$, calculated $f_H$, $C_3$ pole and worst-case offset. Schematic.}

\subsubsection{Frequency response}
\todo[inline]{Bode plot: (a) without $C_3$, (b) with $C_3 = \qty{1}{nF}$, (c) cut-off frequencies marked and compared with the calculation.}

\subsubsection{SPICE analysis}
\begin{figure}[ht]
  \centering
  \includegraphics[width=\linewidth]{sim_noninverting_ac}
  \caption{Non-inverting amplifier without $C_3$, AC analysis (ngspice, TL061 model)}
  \label{fig:sim_noninverting}
\end{figure}
\begin{figure}[ht]
  \centering
  \includegraphics[width=\linewidth]{sim_noninverting_c3_ac}
  \caption{Non-inverting amplifier with $C_3$, AC analysis (ngspice, TL061 model)}
  \label{fig:sim_noninverting_c3}
\end{figure}
\todo[inline]{Compare the simulated cut-off frequencies (\autoref{fig:sim_noninverting} and~\ref{fig:sim_noninverting_c3}) with the calculated ones and the measurement.}

\subsubsection{Offset}
\todo[inline]{Measured output offset, input-referred value, comparison with the data sheet.}

\subsubsection{Slew rate}
\todo[inline]{Oscilloscope screenshot, measured slew rate, comparison with the data sheet.}

\subsection{Integrator}\label{sec:integrator}
\todo[inline]{Design: input square wave $U_{in}$, $f$, target triangle $U_{out}$; $R_1$, $C_3$, $R_4$, $R_5$. Schematic.}

\subsubsection{Frequency response}
\todo[inline]{Bode plot with the band of correct integration marked (\qty{-20}{dB} per decade, above $1/(2\pi R_4 C_3)$).}

\subsubsection{Waveforms}
\todo[inline]{Oscilloscope screenshot: input square wave and output triangle wave.}

\subsubsection{SPICE analysis}
\begin{figure}[ht]
  \centering
  \includegraphics[width=\linewidth]{sim_integrator_ac}
  \caption{Integrator, AC analysis; shaded: band of correct integration (ngspice)}
  \label{fig:sim_integrator_ac}
\end{figure}
\begin{figure}[ht]
  \centering
  \includegraphics[width=\linewidth]{sim_integrator_tran}
  \caption{Integrator, transient analysis with the square-wave input (ngspice)}
  \label{fig:sim_integrator_tran}
\end{figure}
\todo[inline]{Compare the simulated triangle amplitude and band with the design and the measurement.}

\subsection{Differentiator}\label{sec:differentiator}
\todo[inline]{Design: input triangle wave $U_{in}$, $f$, target square $U_{out}$; $R_4$, $C_1$, $R_1$ from the stability condition, $R_5 = R_4$. Schematic.}

\subsubsection{Frequency response}
\todo[inline]{Bode plot with the band of correct differentiation marked (\qty{+20}{dB} per decade, up to $1/(2\pi R_1 C_1)$).}

\subsubsection{Waveforms}
\todo[inline]{Oscilloscope screenshot: input triangle wave and output square wave.}

\subsubsection{SPICE analysis}
\begin{figure}[ht]
  \centering
  \includegraphics[width=\linewidth]{sim_differentiator_ac}
  \caption{Differentiator, AC analysis; shaded: band of correct differentiation (ngspice)}
  \label{fig:sim_differentiator_ac}
\end{figure}
\begin{figure}[ht]
  \centering
  \includegraphics[width=\linewidth]{sim_differentiator_tran}
  \caption{Differentiator, transient analysis with the triangle-wave input (ngspice)}
  \label{fig:sim_differentiator_tran}
\end{figure}
\todo[inline]{Compare the simulated square amplitude and band with the design and the measurement.}
